% % Prevent accidental compilation
% \documentclass{beamer}
% \usetheme{Madrid}
% \usecolortheme{default}

% \title{Thesis Proposal}
% \subtitle{Network Analysis of Venture Capital and Accelerator Ecosystems}
% \author{João Melga}
% \institute{Master IREN}
% \date{June 6, 2025}

% \begin{document}

% \frame{\titlepage}

% \begin{frame}
% \frametitle{Outline}
% \tableofcontents
% \end{frame}

% \section{Motivation}

% \begin{frame}
% \frametitle{Motivation}
% \begin{itemize}
%     \item Venture capital (VC) and accelerator ecosystems are crucial for startup financing
%     \item Complex network relationships between VCs, accelerators, and startups
%     \item Limited understanding of structural patterns and matching mechanisms
%     \item Need for quantitative analysis of ecosystem dynamics
% \end{itemize}
% \end{frame}

% \section{Problem Statement}

% \begin{frame}
% \frametitle{Problem Statement}
% \begin{block}{Macro Research Question}
% How do network structures emerge between venture capitalists and accelerators, and what patterns govern their investment relationships?
% \end{block}

% \begin{itemize}
%     \item Current literature focuses on individual actor behavior
%     \item Gap in understanding ecosystem-level network structures
%     \item Limited application of network theory to VC-accelerator relationships
% \end{itemize}
% \end{frame}

% \section{Literature Review}

% \begin{frame}
% \frametitle{Literature Review}
% \textbf{Relevant Topics}
% \begin{itemize}
%     \item \textbf{Screening}: Focus on entrepreneurial ecosystem
%     \item \textbf{Network Theory}: Structural analysis of relationships
%     \item \textbf{Bipartite Networks}: Two-mode networks with distinct node types
%     \item \textbf{Assortative Matching}: Tendency for similar entities to connect
%     \item \textbf{Nestedness}: Hierarchical interaction patterns
%     \item \textbf{Complexity Metrics}: Network complexity indices
% \end{itemize}
% \vspace{0.5cm}
% \textbf{Existing Gaps}
% \begin{itemize}
%     \item Limited analysis of VC-accelerator intermediary networks
%     \item Lack of complexity measures in ecosystem analysis
% \end{itemize}
% \end{frame}

% \section{Methodology}

% \begin{frame}
% \frametitle{Methodology - Potential Research Directions}
% \begin{columns}
% \begin{column}{0.5\textwidth}
% \textbf{Path 1: Assortative Matching}
% \begin{itemize}
%     \item[\textcolor{red}{\times}] Complex theoretical requirements
%     \item[\textcolor{red}{\times}] Limited data availability
%     \item[\textcolor{red}{\times}] High methodological barriers
%     \item[\textcolor{orange}{?}] Uncertain feasibility
% \end{itemize}
% \end{column}
% \begin{column}{0.5\textwidth}
% \textbf{Path 2: Bipartite Network Analysis}
% \begin{itemize}
%     \item[\textcolor{green}{\checkmark}] Available data
%     \item[\textcolor{green}{\checkmark}] Established methods
%     \item[\textcolor{green}{\checkmark}] Novel application
%     \item[\textcolor{green}{\checkmark}] Original contribution
% \end{itemize}
% \end{column}
% \end{columns}
% \end{frame}

% \begin{frame}
% \frametitle{Methodology - Potential Research Directions}
% \textbf{Proposed Path - Bipartite Network Approach}
% \begin{itemize}
%     \item \textbf{Network Construction}: VCs and accelerators as distinct node types
%     \item \textbf{Edge Definition}: Co-investment relationships
%         \item \textbf{Structural Analysis}: Clustering, centrality measures
%     \item \textbf{Nestedness Calculation}: Hierarchical organization patterns
%     \item \textbf{Complexity Metrics}: Network complexity indices
%     \end{itemize}
% \vspace{0.5cm}
% \textbf{Preliminary Analysis Results (Ex2)}
% \begin{itemize}
%     \item Clustering analysis reveals distinct investment communities
%     \item Evidence of hierarchical structure in VC-accelerator relationships
%     \end{itemize} % Ensure all environments are properly closed
% \end{frame}

% \section{Expected Contributions}

% \begin{frame}
% \frametitle{Expected Contributions}
% \textbf{Theoretical}
% \begin{itemize}
%     \item Novel application of bipartite network theory to VC ecosystems
%     \item Framework for understanding intermediary relationships
%     \item Extension of nestedness concepts to financial networks
%     \item Integration of complexity theory with ecosystem analysis
% \end{itemize} % Ensure all environments are properly closed
% \vspace{0.5cm}
% \textbf{Practical}
% \begin{itemize}
%     \item Insights for VC investment strategy optimization
%     \item Accelerator positioning and partnership strategies
% \end{itemize}
% \end{frame}

% \section{Timeline and Next Steps}

% \begin{frame}
% \frametitle{Timeline and Next Steps}
% \begin{itemize}
%     \item \textbf{Weeks 1-2}: Data collection and preprocessing
%     \item \textbf{Weeks 3-5}: Network construction and basic analysis
%     \item \textbf{Weeks 6-7}: Nestedness and complexity calculations
%     \item \textbf{Weeks 8-10}: Results interpretation and writing
%     \item \textbf{Weeks 11-12}: Thesis completion and defense preparation
% \end{itemize}
% \end{frame}

% \section{Discussion}

% % Review if this part is really necessary
% \begin{frame}
% \frametitle{Questions for Discussion}
% \begin{enumerate}
%     \item Is the bipartite network approach the most promising direction?
%     \item What additional data sources should be considered?
%     \item How can we ensure the originality of complexity measures?
%     \item What are the key theoretical frameworks to incorporate?
%     \item How can we validate our network findings?
% \end{enumerate}
% \end{frame}

% \begin{frame}
% \frametitle{Thank You}
% \begin{center}
% \Large Questions and Discussion
% \end{center}
% \end{frame}

% \end{document}